\documentclass[11pt,a4paper]{article}
\usepackage[T1]{fontenc}
\usepackage{lmodern,microtype}
\usepackage[margin=24mm,headheight=15pt]{geometry}
\usepackage{amsmath,amssymb,booktabs,tabularx,array}
\usepackage{enumitem,xcolor,fancyhdr,xurl}
\usepackage[colorlinks=true,linkcolor=blue!50!black,urlcolor=blue!50!black]{hyperref}

\definecolor{ink}{HTML}{182737}
\definecolor{accent}{HTML}{225878}
\definecolor{pale}{HTML}{EFF5F7}
\color{ink}
\DeclareUnicodeCharacter{2082}{\textsubscript{2}}
\setlength{\emergencystretch}{4em}
\setlength{\parindent}{0pt}
\setlength{\parskip}{5.5pt}
\setlist{nosep,leftmargin=1.55em}
\renewcommand{\arraystretch}{1.25}
\pagestyle{fancy}
\fancyhf{}
\fancyhead[L]{\small\textcolor{accent}{RIEMANN HYPOTHESIS / LEAN SOURCE GUIDE}}
\fancyfoot[L]{\small\textcolor{accent}{Independent technical reading notes}}
\fancyfoot[R]{\small\thepage}
\renewcommand{\headrulewidth}{0.35pt}
\newcommand{\repo}{https://github.com/LiRunGuo/riemann-lean-verification}
\newcommand{\code}[1]{\texttt{#1}}
\newcommand{\src}[2]{\href{\repo/blob/main/lean/src/#1}{\mbox{\nolinkurl{#2}}}}
\newcommand{\evid}[2]{\href{\repo/blob/main/#1}{#2}}
\newcommand{\cb}{Development/AutomorphicArithmetic/CharacterBoxes}
\newcommand{\pc}{Development/AutomorphicArithmetic/PrimitiveCosets}
\newcommand{\lt}{Development/AutomorphicArithmetic/DedekindLFunctions/LandauTheorems}
\newcommand{\tc}{Development/AutomorphicArithmetic/ThetaCorrespondence/ThetaPairLifts/ThirdPlaceSeries}
\newcommand{\hc}{Formalism/SharpRestriction/RamificationFrobenius/TotientGaloisClasses}
\newcommand{\dt}{Development/AutomorphicArithmetic/AutomorphicVarieties/ModularityArgument/ThresholdArguments/DenseThresholdCases}
\newcommand{\RePart}{\operatorname{Re}}
\newcommand{\note}[1]{\par\noindent\colorbox{pale}{\parbox{\dimexpr\linewidth-2\fboxsep\relax}{\small\strut #1\strut}}\par}
\hypersetup{pdftitle={A Technical Guide to the Lean Riemann Hypothesis Submission},
pdfsubject={Theorem statements, proof dependencies, and verification evidence},
pdfkeywords={Riemann hypothesis, Dirichlet L-functions, Lean, formal verification}}

\begin{document}
\thispagestyle{plain}
{\large\color{accent} Technical guide \hfill 20 November 2026}\par
\vspace{8mm}
{\Huge\bfseries The Riemann hypothesis\par}
\vspace{2mm}
{\LARGE A guide to the submitted Lean proof\par}
\vspace{5mm}

This guide follows the Lean project linked in Levent Alpöge's November 2026
announcement. It identifies the precise theorem statements, reconstructs the
analytic deduction, and traces the arithmetic inputs that would establish the
claimed zero-free region. Linked filenames refer to the archived Lean sources
in the verification repository. \href{https://alpo.ge/remark.pdf}{Original announcement}

\note{\textbf{Verification status at drafting.} The targeted verification run checked
seven analytic lemmas using only the three standard Lean axioms. The full build
of the submission's \code{Main} module was still running. The arithmetic
inputs to the central estimate remain open in this audit.
\evid{evidence/status.md}{Live verification record}}

\section{What the submission asserts}

\src{Main.lean}{Main.lean} exposes five declarations. The main RH declaration
has Mathlib's own type \code{RiemannHypothesis}. The Dirichlet assertions
use Mathlib's \code{DirichletCharacter.LFunction} at every positive modulus.

\begin{tabularx}{\linewidth}{@{}>{\raggedright\arraybackslash}p{0.46\linewidth}X@{}}
\toprule
Lean declaration & Mathematical content \\
\midrule
\nolinkurl{riemannHypothesis} & Mathlib's statement of RH. \\
\nolinkurl{riemannZeta_eq_zero_re} & The same zero condition written out in the submission. \\
\nolinkurl{LFunction_ne_zero_of_half_lt_re} & Nonvanishing to the right of the critical line, with the principal character at $s=1$ excepted. \\
\nolinkurl{LFunction_eq_zero_re} & Every zero of every Dirichlet $L$-function in $0<\RePart s<1$ lies on $\RePart s=\tfrac12$. \\
\nolinkurl{LFunction_eq_zero_re_of_isPrimitive} & The critical-strip assertion restricted to primitive characters. \\
\bottomrule
\end{tabularx}

See \src{Statements.lean}{Statements.lean} and
\href{https://github.com/leanprover-community/mathlib4/blob/0df444a360eaa60ab8c11dca51a86af692955474/Mathlib/NumberTheory/LSeries/RiemannZeta.lean\#L185}{the pinned Mathlib definition of RH}.
The term GRH in this guide refers to Dirichlet $L$-functions.

\vspace{2mm}
\textbf{Proof in one line.}
Construct arithmetic families with controlled cohomology and divisibility,
obtain a square-root-scale bound for smoothed prime sums on power intervals,
apply a Landau-type zero-free criterion, and reflect zeros by the functional
equations.

\clearpage
\section{The reduction to right-half-plane nonvanishing}

The proof routes RH and the critical-strip assertions through
\code{UpperTheory.EveryCharacter}. In mathematical notation this proposition is
\[
  \forall N\ge 1,\ \forall \chi\pmod N,\ \forall s\in\mathbb C,
  \qquad
  \RePart s>\tfrac12,\ (\chi,s)\ne(1,1)
  \quad\Longrightarrow\quad L(s,\chi)\ne0.
\]
The source uses the disjunction $\chi\ne1\lor s\ne1$. The exception handles
the pole under the conventions of the formalized $L$-function. The equivalent
right-half-plane input in the final wrapper is \code{Claim.HalfPlane}.
\src{\cb/EveryCharacter.lean}{EveryCharacter.lean}
\src{\cb/HalfPlane.lean}{HalfPlane.lean}

\subsection*{Primitive and imprimitive characters}

For a character $\chi$ modulo $N$, let $\chi_*$ be its primitive character of
conductor $f\mid N$. Passing from $\chi_*$ to $\chi$ introduces a finite product
of Euler factors of the form $1-\chi_*(p)p^{-s}$. If $\RePart s>0$, then
\[
  \bigl|\chi_*(p)p^{-s}\bigr|
     \le p^{-\RePart s}<1 .
\]
No such factor vanishes there. Thus a zero of the imprimitive $L$-function in
the open strip gives a zero for the primitive inducing character. The same
factorization extends a primitive right-half-plane result to all characters.
\src{\cb/ClaimStrip.lean}{ClaimStrip.lean}
\src{\cb/Imprimitive.lean}{Imprimitive.lean}

For primitive $\psi$, the completed functional equation takes a zero at $s$
to a zero of the completed function for $\psi^{-1}$ at $1-s$. In the open
strip, the gamma factor at $s$ is nonzero, so the original $L$-zero gives a
completed zero. At the reflected point the source rewrites the $L$-function
as the completed function divided by its gamma factor. If
$0<\RePart s<\tfrac12$, then $\RePart(1-s)>\tfrac12$, and the reflected point
contradicts right-half-plane nonvanishing. Zeros with $\RePart s>\tfrac12$
contradict it directly. The critical line is the remaining case.
\src{\cb/ClaimStrip.lean}{ClaimStrip.lean}

\subsection*{The zeta function}

The character of modulus $1$ satisfies
$L(s,\chi)=\zeta(s)$ in Mathlib. The route through
\code{Formalism.VW.rh\_of\_every} uses this equality and the zeta functional
equation. Mathlib's RH statement explicitly excludes the negative even
integers and $s=1$. The source also uses $\zeta(0)\ne0$. An alternative
reduction, \code{Claim.zeta}, treats the nonpositive half-plane via the
completed zeta function and the established nonvanishing on $\RePart s\ge1$.
The reduction audit checked the analytic lemmas underlying both routes.
\src{Supplements/Kernels/MainTheorems.lean}{MainTheorems.lean}
\src{Formalism/GammaFactorNonvanishing/CritLine.lean}{CritLine.lean}
\src{Formalism/GammaFactorNonvanishing/RhOfRight.lean}{RhOfRight.lean}

\clearpage
\section{The analytic criterion on power intervals}

\subsection*{Smoothing and normalization}

Write $\Lambda(n)$ for the von Mangoldt function. The zeta component is the
centered sum
\begin{align*}
 S_\zeta(x)&=\sum_{n\le x}\Lambda(n)(x/n-1),\\
 G_\zeta(x)&=S_\zeta(x)-x\log x+(1+\gamma)x,
\end{align*}
where $\gamma$ is the Euler--Mascheroni constant. For a nonprincipal character
the complex version of the centered sum replaces $\Lambda(n)$ by
$\Lambda(n)\chi(n)$ and adds $xL'(1,\chi)/L(1,\chi)$. The source takes real
parts and combines a finite family with positive weights. The cyclic level of
a primitive character includes the primitive characters induced by its
powers and includes the trivial character. The resulting level sum is
\code{Level.Gm}. \code{level\_family} identifies that sum with the
weighted family in the analytic criterion, preserving the zero-free goal.
\src{Formalism/StatementOC.lean}{StatementOC.lean}
\src{\hc/Analytic.lean}{Analytic.lean}
\src{\hc/ExplicitFormulaErrorOfHostConditions.lean}{ExplicitFormulaErrorOfHostConditions.lean}
\src{\pc/LevelFamily.lean}{LevelFamily.lean}

\subsection*{Why zeros leave uncancelled poles}

For $\RePart s>1$ the smoothing factor satisfies
\[
 \int_n^\infty (x/n-1)x^{-s-1}\,dx=\frac{n^{-s}}{s(s-1)}.
\]
The Dirichlet series for the smoothed sum therefore enters through
$-L'(s,\chi)/L(s,\chi)$. At a zero $\rho$ of multiplicity $m$, that logarithmic
derivative contributes $-m/(s-\rho)$ in its principal part. Centering removes
the terms at $s=1$. The residues at a zero of the family take the form
\[
                  -\frac{S}{\rho(\rho-1)},\qquad S>0 ,
\]
where $S$ is a positive weighted sum of multiplicities, including the
conjugate-character terms used to express real parts. Since all contributions
at a fixed $\rho$ have the same complex factor, they cannot cancel. In
particular the residue is positive at a real zero in $(0,1)$. This sign
matters because the argument uses an upper bound on a real-valued sum.
\src{\lt/poleData.lean}{poleData.lean}

\subsection*{From a window bound to nonvanishing}

Let $0<\theta<1$ and $\delta>0$. The criterion assumes that for every lower
threshold there exists a larger $Y$ such that $G_m(y)\le C y^\theta$ on
the whole interval $[Y,Y^{1+\delta}]$. After $u=\log y$, high derivatives of
the Laplace transform insert the kernel $u^k\exp(-\sigma u)/k!$.
The proof chooses $k$ so this kernel concentrates inside
$[\log Y,(1+\delta)\log Y]$ and bounds the tails. Cauchy estimates compare
the derivative bound with the growth forced by a nearby pole. The argument
handles the real-pole case using the positive residue. The conclusion is
nonvanishing for the family on $\RePart s>\theta$, with the region
$\RePart s\ge1$ handled by existing theorems.
\src{\lt/landauTails.lean}{landauTails.lean}
\src{\lt/landauWindow.lean}{landauWindow.lean}
\src{\lt/ZeroFreeOfRpowWindow.lean}{ZeroFreeOfRpowWindow.lean}

\clearpage
\section{What the arithmetic argument must supply}

For the cyclic level $\mathcal L$ of each primitive character, the argument
requires families of windows on which
\[
  G_{\mathcal L}(x)
     \le x^{1/2+(e_D+e_L+7)\varepsilon}.
\]
The final assembly uses the power $\lambda=3/2$ for the windows. The constants
$e_D,e_L$ are chosen before the quantifier over all sufficiently small
$\varepsilon>0$. For a candidate zero with real part $\beta>1/2$, the code
chooses a $\theta$ strictly between $1/2$ and $\min(\beta,1)$ and makes
$(e_D+e_L+7)\varepsilon$ small enough to apply the analytic criterion.
The thresholds and the family may depend on $\varepsilon$.
A bound at isolated points would leave the window hypothesis unmet.
\src{\hc/TopWin.lean}{TopWin.lean}
\src{\cb/ZeroFreeAllOfTopG.lean}{ZeroFreeAllOfTopG.lean}
\src{\tc/TopPlug.lean}{TopPlug.lean}

\subsection*{The bridge from arithmetic data to the bound}

The theorem that produces the quantitative bound is \code{UpperTheory.Comparison₂.core\_gm}.
It takes two inputs at scale $x$.

\begin{enumerate}
\item \code{EAtH} relates a level's primitive characters to data on an
auxiliary arithmetic host and controls comparison errors.
\item \code{CoreAtH} records a witness with the identities, size bounds,
and error estimates needed by the prime-sum argument.
\end{enumerate}

The proof matches the weighted sum to the host expression, applies the
chain estimate \code{Deduction.ch\_top} and the comparison
\code{Deduction.ch\_core}, and absorbs fixed factors into the spare powers
of $x^\varepsilon$. The family-level theorem chooses a suitable member
for each $x$ in a window. The analytic theorem then consumes the resulting
inequality and the window quantifiers.
\par{\small\raggedright \src{\pc/CoreAtH.lean}{CoreAtH.lean}
\src{\pc/CoreGm.lean}{CoreGm.lean}
\src{\pc/WindowInequalityOfArtinInputs.lean}{WindowInequalityOfArtinInputs.lean}
\src{\pc/LevelGmLeOnFamilyScales.lean}{LevelGmLeOnFamilyScales.lean}\par}

\subsection*{How the family requirements are assembled}

\code{zeroFreeAll\_of\_topG} separates the existence of hosts from the
construction of compatible estimates over a family. The supply argument
\code{supplyChi} combines a cohomology window for a primitive character
with admissible arithmetic classes. The other input, \code{TopG}, produces
the family-level estimates under that supply hypothesis and the relevant
cohomology bounds. \code{topDesWin} collects the design-existence,
divisibility, rank and global compatibility hypotheses used to obtain
\code{ChainTied} witnesses on the windows. These witnesses feed
\code{CoreAtH}. This is a useful boundary for reviewing whether each
downstream estimate applies to the hosts that are actually supplied.
\src{\cb/SupplyChi.lean}{SupplyChi.lean}
\src{\cb/HostOfField.lean}{HostOfField.lean}
\src{\cb/DesWin.lean}{DesWin.lean}

\clearpage
\section{The two arithmetic branches}

\subsection*{Auxiliary fields and the cohomology window}

\code{IsChiField} describes the auxiliary field through an embedding into
a cyclotomic field and a subgroup condition involving the given character
and two quadratic conditions. \code{WinPair} prescribes the relative sizes
and congruence conditions on auxiliary primes $q,d$. The field bound
\code{VVle} demands finiteness of the relevant lengths and bounds the
weighted sums over finite prime sets, with weights $\log\ell$.
Here the length is that of a localization kernel in degree-one Galois
cohomology for the twist by $-e$. A nontrivial point for review is that
the field and local data actually constructed for the prime-sum estimate
meet the conditions of this bound.
\src{\cb/ExistsPageN.lean}{ExistsPageN.lean}
\src{Formalism/GaloisCohomology/Statement.lean}{Statement.lean}
\src{Development/AutomorphicArithmetic/ShortNorms/supplyOfAdmOf.lean}{supplyOfAdmOf.lean}

The proposition \code{VWinAt} supplies a window for each sufficiently small
$\varepsilon$ and admissible choice of $e(x)$. Its pointwise target
\code{VAt} uses a cutoff $\lfloor x^{401}\rfloor$ in the specialization
\code{VWin}, with cohomology bound $\exp(x^\varepsilon/2)$.
The route through \code{vwin\_of\_row} combines local--global compatibility,
base change, automorphy inputs and several counting bounds. The proof
packages these inputs before applying \code{vwin\_of\_counts}.
\src{\dt/VwinOfRow.lean}{VwinOfRow.lean}
\src{\dt/WindowBoundOfLiteratureLevel.lean}{WindowBoundOfLiteratureLevel.lean}
\src{Development/AutomorphicArithmetic/AutomorphicVarieties/HeightAndVanishing/JointWinOfCounts.lean}{JointWinOfCounts.lean}

\subsection*{Theta matching, covolumes and divisibility}

The other branch constructs the hypotheses used for the prime-sum
comparison. The assembly passes theta-integral agreement and local matching,
a covolume bound, a local Langlands existence result, and the cohomology
window to \code{every\_character\_of\_theta\_integrals\_local\_matching}.
Inside that route the supplied results are used to construct the required
nonempty family hypotheses. At the design stage,
\code{divinN\_of\_hypV} turns a cohomology input into divisibility data
using additional link and rank hypotheses. The window bound depends on
this compatibility as well as on analytic error estimates.
\par{\small\raggedright \src{\tc/EveryCharacterOfLiteratureOfVWin.lean}{EveryCharacterOfLiteratureOfVWin.lean}
\src{\tc/SiegelEisensteinSecondTermHypotheses.lean}{SiegelEisensteinSecondTermHypotheses.lean}
\src{\cb/DivinNOfHypV.lean}{DivinNOfHypV.lean}\par}

\note{\textbf{Where to start a deeper review.} Check the proof of the
local--global compatibility row in the form used by \code{vwin\_of\_row},
then follow one concrete character all the way from \code{VWinAt} to its
\code{CoreAtH} witness. This tests whether the arithmetic constructions
supply the same data and uniformity that the analytic comparison consumes.}

\clearpage
\section{Reading order and verification evidence}

Start from the actual proof terms in \src{Main.lean}{Main.lean}.
The principal route passes through \code{allChi} in
\src{Supplements/Kernels/MainTheorems.lean}{MainTheorems.lean} to
\code{UpperTheory.FinalTheorems.everyCharacter} in
\src{\cb/MainConclusion.lean}{MainConclusion.lean}.
The final declaration supplies \code{galoisLocalGlobalStrong\_recParameter}
from the theta-lift and Galois-representation development. Follow the
named declarations rather than treating every file imported by a module as
a direct mathematical premise.
\src{Development/AutomorphicArithmetic/Foundations/GaloisRepresentations/IntegralStructures/RowSrc.lean}{RowSrc.lean}

\begin{tabularx}{\linewidth}{@{}>{\raggedright\arraybackslash}p{0.26\linewidth}X@{}}
\toprule
Question & Files to open \\
\midrule
What is being proved? &
\src{Main.lean}{Main.lean},
\src{Statements.lean}{Statements.lean} \\
How is the upper bound used? &
\src{\lt/ZeroFreeOfRpowWindow.lean}{ZeroFreeOfRpowWindow.lean},
\src{\hc/TopWin.lean}{TopWin.lean} \\
Where is the upper bound assembled? &
\src{\pc/CoreGm.lean}{CoreGm.lean},
\src{\cb/ZeroFreeAllOfTopG.lean}{ZeroFreeAllOfTopG.lean} \\
Where are both branches joined? &
\src{\cb/MainConclusion.lean}{MainConclusion.lean},
\src{\tc/EveryCharacterOfLiteratureOfVWin.lean}{EveryCharacterOfLiteratureOfVWin.lean} \\
\bottomrule
\end{tabularx}

\subsection*{What has been checked}

{\raggedright
The isolated reduction job \code{157173} completed with exit status
zero. Lean checked the two conditional deductions from
\code{Claim.HalfPlane} to RH and Dirichlet GRH. The axiom audit covered
\code{Claim.prim\_reflect}, \code{Claim.strip}, \code{Claim.zeta},
\code{Claim.factor\_ne\_zero}, \code{Formalism.VW.rh\_of\_right},
\code{Deduction.zeroFree\_of\_rpow\_window} and
\code{UpperTheory.Comparison₂.formB}. Each reported
\code{propext}, \code{Classical.choice} and \code{Quot.sound}.
The full job \code{157166} was still compiling the import closure of
\code{Main} when these notes were drafted. Its recorded procedure checks the
top-level statements and transitive axioms and then runs
\code{leanchecker --fresh Main}.
\evid{evidence/reduction-audit.txt}{Reduction job output}
\evid{verification/AuditReduction.lean}{Reduction audit source}
\evid{evidence/status.md}{Full job status}\par}

\subsection*{Reproducibility}

The public archive consists of fourteen concatenated segments. The public
verification repository preserves 175,187 submitted Lean source files and
the three project configuration files. The toolchain is Lean \code{v4.33.1},
and the Mathlib pin is
\nolinkurl{0df444a360eaa60ab8c11dca51a86af692955474}.
\evid{evidence/archive.sha256}{Segment checksums}
\evid{verification/build-on-slurm.sh}{Full build script}

The evidence supports the analytic reductions under their stated inputs.
The decisive pending check is whether the submitted arithmetic constructions
supply the window estimate for every required character.

\end{document}
